% Generated draft: requires mathematical and source review.
\documentclass[11pt]{article}
\usepackage[margin=25mm]{geometry}
\usepackage{fontspec}
\setmainfont{Latin Modern Roman}
\usepackage{amsmath,amssymb,mathtools,unicode-math}
\setmathfont{Latin Modern Math}
\usepackage{enumitem,needspace,xurl,hyperref}
\setlength{\parindent}{0pt}
\setlength{\parskip}{5pt}
\setlength{\emergencystretch}{3em}
\newcommand{\N}{\mathbb{N}}
\newcommand{\Z}{\mathbb{Z}}
\newcommand{\Q}{\mathbb{Q}}
\newcommand{\R}{\mathbb{R}}
\newcommand{\C}{\mathbb{C}}
\newcommand{\F}{\mathbb{F}}
\newcommand{\A}{\mathbb{A}}
\newcommand{\PP}{\mathbb{P}}
\newcommand{\E}{\mathbb{E}}
\newcommand{\eps}{\varepsilon}
\newcommand{\dd}{\,\mathrm{d}}
\begin{document}
\section*{UnsolvedMath ID: 30002161}
\textbf{Source problem code:} catalogue entry
\subsection{Definitions and mathematical statement}
Write \(E[n]=\ker([n]\colon E\to E)\) for the \(n\)-torsion group scheme of an elliptic curve \(E/k\), whenever the torsion index \(n\) and the cohomological setting have been specified. The supplied record gives a restriction map in a Selmer model,
\[
\beta_E\colon H^1(k,E[n])\longrightarrow H^1(A,E[n]),
\qquad D_E=\operatorname{im}(\beta_E).
\]
The record does not specify the allowed fields \(k\), the allowed values or quantification of \(n\), the meaning of \(A\), or the cohomology theory and local conditions used to define this map. These symbols are therefore retained as source notation, rather than assigned additional hypotheses.

For an abelian group \(B\), a subgroup \(D\subseteq B\) is an algebraic direct summand if there exists a subgroup \(C\subseteq B\) such that \(B=D\oplus C\). On this algebraic interpretation, the requested property is
\[
\exists C_E\subseteq H^1(A,E[n])\text{ a subgroup such that }
H^1(A,E[n])=D_E\oplus C_E.
\]
The source does not say whether the intended splitting must additionally respect a topology or other structure.

\textbf{Question, with unresolved parameters retained.} Is \(D_E\) a direct summand of \(H^1(A,E[n])\) for \(100\%\) of elliptic curves \(E\) over \(k\)? To express the density quantifier precisely, one would need the intended finite counting families \(\mathcal F(T)\) of elliptic curves over \(k\), including the equivalence relation and ordering used. Once these are specified, the algebraic version of the question has the form
\[
\lim_{T\to\infty}
\frac{\#\{E\in\mathcal F(T):D_E\text{ is an algebraic direct summand of }H^1(A,E[n])\}}
{\#\mathcal F(T)}=1,
\]
with the denominator required to be nonzero eventually. This is a conditional formulation of the supplied density-one question; the record does not determine a particular counting family or the quantifiers on \(k\) and \(n\).

The supplied background mentions a known rational-torsion case and the Poonen--Rains model as insufficient to settle the general question. It supplies no precise hypotheses or theorem statement for that special case, so neither is incorporated as an additional mathematical assertion.
\subsection{Short English statement}
The question asks whether the image of the Selmer restriction map splits off its target for asymptotically \(100\%\) of elliptic curves over \(k\). A fully specified statement still requires the field and torsion-index hypotheses, the meaning of the target cohomology, the type of splitting, and the counting convention. The supplied open status is dated curation, not an independently verified conclusion.
\subsection{Sources}
\begin{itemize}
\item Supplied local UnsolvedMath registry record, permanent ID \(30002161\), title ``Direct-Summand Property for Selmer Restriction Images,'' source problem code \texttt{catalogue entry}, in the Oberwolfach Reports Open Problems collection. No individual proposer is supplied; the proposed year is recorded as 2012.
\item Supplied report citation: \emph{Arithmetic Geometry (2013)}, Oberwolfach Reports, pp.\ 2335--2388, DOI \texttt{10.4171/owr/2012/38}. Source URL and literature-source link labelled ``OWR report 38'': \url{https://doi.org/10.4171/owr/2012/38}. The supplied background attributes the problem to workshop open-problem discussions documented in this report.
\item Supplied literature-source link labelled ``Poonen--Rains'': \url{https://doi.org/10.1090/S0894-0347-2012-00709-9}. No fuller bibliographic entry or precise supporting theorem is supplied.
\item Literature assessment dated exactly \texttt{2026-08-21}: the supplied curation labels the problem open and states that the known rational-torsion case and the Poonen--Rains model do not settle the general direct-summand question. This is supplied dated curation, not independently verified literature assessment.
\item The registry labels the statement \texttt{corrected\_verified} and says that the official report poses the same direct-summand question with a density-one quantifier. These are supplied verification claims; no source was independently visited or checked for this draft.
\end{itemize}
\end{document}
